\documentclass[10pt]{article}
\providecommand{\FIGDIR}{fig}
\newcommand{\DIGESTDATE}{2 September 2026}
% ===================== PM Research Watch - shared preamble =====================
\usepackage[T1]{fontenc}
\usepackage[utf8]{inputenc}
\usepackage{amsmath,amssymb}
\usepackage{tgpagella}
\usepackage{tgheros}
\renewcommand{\sfdefault}{qhv}
\renewcommand{\ttdefault}{lmtt}
\usepackage{microtype}
\usepackage[english]{babel}

\usepackage[a4paper,top=17mm,bottom=17mm,left=17mm,right=17mm,
            headsep=6mm,footskip=10mm]{geometry}
\usepackage{graphicx}
\usepackage{xcolor}
\usepackage{colortbl}
\usepackage{tikz}
\usetikzlibrary{calc}
\usepackage[most]{tcolorbox}
\usepackage{titlesec}
\usepackage{fancyhdr}
\usepackage{booktabs}
\usepackage{array}
\usepackage{longtable}
\usepackage{enumitem}
\usepackage{multicol}
\usepackage{ragged2e}
\usepackage[hidelinks]{hyperref}

% ------------------------------------------------------------------ palette
\definecolor{Ink}    {HTML}{14262E}
\definecolor{Deep}   {HTML}{0F4C5C}
\definecolor{Teal}   {HTML}{1F7A8C}
\definecolor{Sky}    {HTML}{4FA3C4}
\definecolor{Sage}   {HTML}{6FA287}
\definecolor{Amber}  {HTML}{E3A008}
\definecolor{Clay}   {HTML}{B5643C}
\definecolor{Coral}  {HTML}{D1495B}
\definecolor{Violet} {HTML}{6A5B9E}
\definecolor{Slate}  {HTML}{4A6572}
\definecolor{Paper}  {HTML}{FBF9F5}
\definecolor{Mist}   {HTML}{EEEAE2}
\definecolor{Rule}   {HTML}{D9D3C9}

\pagecolor{Paper}
\color{Ink}

\hypersetup{colorlinks=true, linkcolor=Deep, urlcolor=Teal, citecolor=Teal,
            pdftitle={PM Research Watch}, pdfauthor={Automated daily digest}}

% ------------------------------------------------------------------ headings
\titleformat{\section}
  {\sffamily\bfseries\large\color{Deep}}{}{0pt}
  {}[{\vspace{1.2mm}\color{Amber}\titlerule[1.1pt]}]
\titlespacing*{\section}{0pt}{5.5mm}{2.8mm}

\titleformat{\subsection}
  {\sffamily\bfseries\normalsize\color{Teal}}{}{0pt}{}
\titlespacing*{\subsection}{0pt}{3.6mm}{1.4mm}

\setlist[itemize]{leftmargin=4.2mm, itemsep=1.1mm, topsep=1.3mm, parsep=0pt,
                  label=\textcolor{Teal}{\footnotesize$\blacktriangleright$}}

% ------------------------------------------------------------------ header/footer
\pagestyle{fancy}
\fancyhf{}
\renewcommand{\headrulewidth}{0pt}
\renewcommand{\footrulewidth}{0pt}
\fancyhead[L]{\sffamily\scriptsize\color{Slate}PM RESEARCH WATCH\;\textcolor{Rule}{\textbar}\;\DIGESTDATE}
\fancyhead[R]{\sffamily\scriptsize\color{Slate}Particulate matter monitoring \& health effects}
\fancyfoot[C]{\sffamily\scriptsize\color{Slate}\thepage}

% ------------------------------------------------------------------ components
% metric tile
\newcommand{\metric}[3]{%
  \begin{tcolorbox}[enhanced, colback=white, colframe=Rule, boxrule=0.5pt,
    arc=1.6mm, left=2.2mm, right=2.2mm, top=1.8mm, bottom=1.8mm,
    borderline west={1.7mm}{0pt}{#3}, nobeforeafter, valign=center,
    height=17mm, width=\linewidth]
    \RaggedRight
    {\sffamily\bfseries\fontsize{17}{17}\selectfont\color{#3}#1}\\[0.5mm]
    {\sffamily\fontsize{6.5}{7.4}\selectfont\color{Slate}#2}
  \end{tcolorbox}}

% pull-quote / key insight box
\newtcolorbox{keybox}[1][]{enhanced, colback=Mist, colframe=Mist,
  arc=1.6mm, left=3.2mm, right=3.2mm, top=2.4mm, bottom=2.4mm,
  borderline west={1.7mm}{0pt}{Amber}, #1}

% subtopic band
\newcommand{\band}[2]{%
  \vspace{3.2mm}\noindent
  \begin{tcolorbox}[enhanced, colback=#2, colframe=#2, arc=1.2mm,
    left=2.6mm, right=2.6mm, top=1.1mm, bottom=1.1mm, boxrule=0pt,
    width=\linewidth, nobeforeafter]
    {\sffamily\bfseries\small\color{white}#1}
  \end{tcolorbox}\vspace{1.6mm}}

% one paper entry: {short cite}{title}{journal}{doi}{body}
\newcommand{\paperentry}[5]{%
  \par\vspace{1.9mm}\noindent
  {\sffamily\bfseries\footnotesize\color{Ink}#1}\;%
  {\sffamily\fontsize{7}{8}\selectfont\color{Slate}#3}\\[0.4mm]
  {\footnotesize\itshape\color{Deep}#2}\\[0.6mm]
  {\footnotesize #4}\\[0.3mm]
  {\fontsize{6.6}{7.6}\selectfont\color{Slate}\texttt{doi:}\,\href{https://doi.org/#5}{\texttt{\detokenize{#5}}}}
  \par}

% caption style
\newcommand{\figcap}[1]{\par\vspace{1.0mm}
  {\fontsize{7.2}{8.6}\selectfont\color{Slate}\textsf{\textbf{Fig.}}\; #1\par}\vspace{1.5mm}}

\providecommand{\permil}{\textperthousand}

\setlength{\parindent}{0pt}
\setlength{\parskip}{0pt}
\renewcommand{\arraystretch}{1.18}


\begin{document}

% ==================================================================== MASTHEAD
\thispagestyle{empty}
\begin{tikzpicture}[remember picture, overlay]
  \fill[Deep] (current page.north west) rectangle
      ($(current page.north east)+(0,-67mm)$);
  \fill[Amber] ($(current page.north west)+(0,-67mm)$) rectangle
      ($(current page.north east)+(0,-68.4mm)$);
  \foreach \x/\y/\r in {22/12/0.9, 41/26/1.5, 63/9/0.7, 88/31/1.1, 112/17/1.9,
                        138/29/0.8, 159/13/1.3, 178/54/1.0, 196/21/1.6,
                        30/44/0.6, 74/57/1.2, 125/48/0.9, 168/6/0.8, 191/59/0.7,
                        52/18/1.0, 100/60/0.6, 148/20/1.1, 66/49/1.4, 205/45/0.8}
    \fill[Sky, opacity=0.30] ($(current page.north west)+(\x mm,-\y mm)$) circle (\r mm);
\end{tikzpicture}

\vspace*{4mm}
{\sffamily\fontsize{7.6}{9}\selectfont\color{Sky}%
 AUTOMATED DAILY LITERATURE DIGEST\;\textbullet\;PhD RESEARCH WATCH\par}
\vspace{2.2mm}
{\sffamily\bfseries\fontsize{27}{29}\selectfont\color{white}%
 Particulate Matter\par}
{\sffamily\fontsize{27}{29}\selectfont\color{Sky}%
 Monitoring \& Health Effects\par}
\vspace{5.0mm}
{\sffamily\fontsize{8.4}{10}\selectfont\color{white}%
 Issue 2026-09-02\;\;\textcolor{Amber}{\textbar}\;\;Entry window 2 September 2026 (single day)\;\;%
 \textcolor{Amber}{\textbar}\;\;15 records in scope, 62 screened out\par}
\vspace{18mm}

% -------------------------------------------------------------- metric strip
\noindent
\begin{minipage}[t]{0.190\linewidth}\metric{15}{RECORDS IN SCOPE\\(62 SCREENED OUT)}{Deep}\end{minipage}\hfill
\begin{minipage}[t]{0.190\linewidth}\metric{6}{SUBTOPIC CLUSTERS\\OF TEN OCCUPIED}{Teal}\end{minipage}\hfill
\begin{minipage}[t]{0.190\linewidth}\metric{4}{EFFECT ESTIMATES\\(TWO SOURCES)}{Amber}\end{minipage}\hfill
\begin{minipage}[t]{0.190\linewidth}\metric{6}{RECORDS SHIPPING\\ON METADATA ONLY}{Clay}\end{minipage}\hfill
\begin{minipage}[t]{0.190\linewidth}\metric{9}{RECORDS WITH NO\\HEALTH ENDPOINT}{Sky}\end{minipage}

\vspace{6mm}

% -------------------------------------------------------------- signal of day
\section*{\sffamily\bfseries\color{Deep}Signal of the day}
\vspace{-1.5mm}{\color{Amber}\rule{\linewidth}{1.1pt}}\vspace{2.5mm}

\begin{keybox}
{\sffamily\bfseries\small\color{Deep}A national STEMI register and a Mendelian-randomisation
study land on the same day at opposite ends of the effect-size scale --- and the gap
between them is the whole methodological problem of PM epidemiology in one figure.}\\[1.4mm]
\footnotesize
De Potter et al.\ analysed \textbf{every ST-elevation myocardial infarction admission in
Belgium} over nine years ($n$=51{,}544) in a time-stratified case-crossover design with
residential exposure from high-resolution models. Per IQR, same-day PM$_{2.5}$ gave
\textbf{OR 1.022 (1.008--1.035)} --- a 2\% acute excess, tight CI, and the kind of number
a design that self-matches on person and season is built to produce. Li H et al.\ instead
instrument lifetime PM$_{2.5}$ genetically and report \textbf{OR 1.897 (1.063--3.388)}
for childhood asthma, with \textbf{FUT10 mediating 19.3\%} of the effect and 70
TWAS-implicated genes enriched in lysosomal and leukocyte-mediated immunity.
Both are defensible; they are not measuring the same thing. The case-crossover isolates a
trigger effect over hours with the chronic exposure held constant by design; the MR
estimates a lifetime contrast with the acute component invisible. Reading either as
``the'' PM effect is a category error, and any sensor network that claims to serve
epidemiology has to declare which of the two it is resolving --- minutes-to-days for the
first, cumulative annual burden for the second. They imply completely different siting,
duty cycle and calibration requirements.
\end{keybox}

\vspace{3mm}

% -------------------------------------------------------------- exec brief
\section*{\sffamily\bfseries\color{Deep}What changed today}
\vspace{-1.5mm}{\color{Amber}\rule{\linewidth}{1.1pt}}\vspace{2.5mm}

\footnotesize
\begin{itemize}
\item \textbf{Six of fifteen records ship on metadata only} --- the worst abstract
availability this watch has recorded. Elsevier deposited five \textit{Atmos.\ Environ.},
\textit{Environ.\ Pollut.}, \textit{Sci.\ Total Environ.} and \textit{Build.\ Environ.}
records to Crossref with no abstract, none is yet in PubMed or Europe PMC, and Semantic
Scholar returns 404 or a null abstract for all of them. They are counted and listed;
\textbf{no number is attributed to any of them}, and four are on-topic enough
(PM$_{2.5}$-bound metals, solid-fuel PAH biomarkers, heritage-site particle sensing,
landfill receptor monitoring) that they belong on the abstract retry list. All six are
scored \textit{Metadata only} in the architecture figure rather than being given a design
guessed from the title.

\item \textbf{Two effect estimates on the same acute contrast, in the same paper,
disagree by pollutant.} De Potter et al.\ report same-day \textbf{PM$_{2.5}$ OR 1.022
(1.008--1.035)} against \textbf{NO$_2$ OR 1.044 (1.025--1.063)} per IQR, with O$_3$
strongest at lag 0--2 (\textbf{1.063, 1.034--1.090}). NO$_2$ doubling the PM$_{2.5}$
coefficient in a single-pollutant comparison is the traffic-proximity confounding this
literature never resolves --- and in-hospital mortality (11.8\%) was driven not by any
pollutant but by individual and area income and retirement status.

\item \textbf{A filtration material with a number worth writing down.} Lv H et al.'s
single-component polyimide gradient aerogel reaches \textbf{99.996\% PM$_{0.3}$ efficiency
at 45.08\,Pa initial pressure drop} with \textbf{127\,g\,m$^{-2}$ dust-holding capacity},
holding 99.91\% at \textbf{300\,\textdegree C}. The gradient architecture buys depth
filtration rather than surface caking, which is why the dust-holding number is the
interesting one, not the efficiency.

\item \textbf{Personal exposure during exercise, measured rather than modelled.} Engeroff
et al.\ ran 60 instrumented 2\,km runs across a park, a residential street and a major
road at two times of day. Lowest PM across all fractions: \textbf{midday in the park};
highest PM$_{10}$: \textbf{morning on the same park route}. Route choice and rush-hour
avoidance both work, but the meteorological covariates (humidity, pressure, wind,
temperature) carried real signal --- which is a warning about drawing route conclusions
from a single campaign.

\item \textbf{Sulfate aerosol water turns out to be photochemically active on its own.}
Petersen et al.\ show the sulfate anion absorbs surface sunlight and ejects an electron to
form SO$_4^{\bullet-}$ on a sub-microsecond timescale, at rates \textbf{comparable to
measured hydroxyl-radical production in aerosols}. Biomass-fire chromophores photosensitise
the pathway, raising radical formation \textbf{up to five-fold} and more strongly at low
pH. The consequence is not marginal: secondary brown carbon from this chemistry
\textbf{flips the direct radiative forcing of sulfate-containing aerosol from negative to
positive}.

\item \textbf{No sensing-side calibration or drift work today.} \textit{Sensing,
forecasting \& instrumentation} holds two records, both metadata-only. OpenAlex was 403
and arXiv 429 again, Consensus was not re-queried after returning three
already-seen records on the 1 September run, and PubMed's instrumentation axis returned
\textbf{zero} hits for this entry date.
\end{itemize}

\clearpage

% ==================================================================== ANALYTICS
\section*{\sffamily\bfseries\color{Deep}Corpus analytics}
\vspace{-1.5mm}{\color{Amber}\rule{\linewidth}{1.1pt}}\vspace{3mm}

\noindent\includegraphics[width=\linewidth]{\FIGDIR/f1_subtopics.png}
\figcap{Subtopic assignment of the 15 in-scope records. \textbf{Only six of ten clusters
are occupied} --- \textit{Mechanistic toxicology}, \textit{Neuro / mental health},
\textit{Burden, policy \& mitigation} and \textit{Other clinical endpoints} are all empty,
the first day since 24 August with four or more simultaneously vacant.
\textit{Occupational \& indoor} at 4 is unusually large for this watch and is genuinely
heterogeneous: a filtration material, a surgical-smoke bench test, a solid-fuel biomarker
study and an infiltration metric.}

\vspace{2mm}
\noindent\includegraphics[width=\linewidth]{\FIGDIR/f2_design_endpoint.png}
\figcap{Left --- study architecture. \textbf{\textit{Metadata only} is the largest group
at 6 of 15} --- these records are scored as an absence of information rather than assigned
a design inferred from their titles, which is why the real architectures behind them do
not appear anywhere in this panel. Of the nine that do carry a design,
\textit{Modelling / inventory} and \textit{Experimental / toxicology} hold 3 each and the
single \textit{Observational --- acute} entry is the Belgian case-crossover. Right ---
health endpoint. \textbf{Nine of fifteen (60\%)} carry no human health endpoint, below
1 September's 67\%; four of the six that do are metadata-only or MR-derived, so the
observed-incidence content of this issue is one paper.}

\vspace{2mm}
\noindent\includegraphics[width=\linewidth]{\FIGDIR/f3_metric_geography.png}
\figcap{Left --- particle metric. The eight records in \textit{PM (unspeciated) /
emissions} are inflated by the metadata-only entries, whose particle metric can only be
inferred from the title; treat that bar as an artefact of missing abstracts rather than a
statement about the day's measurement practice. Right --- geographic provenance. Europe
leads at 5 --- Belgium and Germany named, three coded only to \textit{Europe} because two
are multi-country and one ships on metadata --- then China 4, North America 2, Latin
America 1. \textbf{Sub-Saharan Africa, South Asia, Southeast Asia
and Oceania are all absent}; Latin America returns for the first time since 28 August, via
a heritage-conservation sensing study.}

\vspace{2mm}
\noindent\includegraphics[width=\linewidth]{\FIGDIR/f6_lifecourse.png}
\figcap{Life-course windows addressed; counts are non-exclusive. \textbf{Puberty and
adolescence is empty}. The childhood-asthma MR is scored against both the in-utero and
early-childhood cells because a genetically instrumented lifetime exposure has no
identifiable window --- which is precisely the interpretive weakness of the design, and
scoring it in a single cell would hide that.}

\vspace{2mm}
\noindent\includegraphics[width=\linewidth]{\FIGDIR/f4_heatmap.png}
\figcap{Subtopic $\times$ study-architecture cross-tabulation. Sparse by construction at
$n$=15, and the \textit{Metadata only} column carries 6 of the 15 --- including
\emph{both} \textit{Sensing, forecasting \& instrumentation} records, which is why that
row contributes no method detail at all. Read this figure as a map of what is missing as
much as of what is present.}

\vspace{2mm}
\noindent\includegraphics[width=\linewidth]{\FIGDIR/f5_forest.png}
\figcap{Four estimates from \textbf{two} sources. The three Belgian case-crossover odds
ratios are per-IQR, same-day (or lag 0--2 for O$_3$) contrasts within one register and are
not independent of one another; the childhood-asthma odds ratio is a Mendelian
randomisation on a lifetime genetically instrumented exposure. They share an axis but not
a scale of inference --- the MR point estimate is roughly forty times larger in excess
risk (89.7\% vs 2.2\%) and its interval spans a factor of three. Do not read the spread as
heterogeneity between studies of the same quantity.}

\clearpage

% ==================================================================== DIGEST
\section*{\sffamily\bfseries\color{Deep}Paper-level digest}
\vspace{-1.5mm}{\color{Amber}\rule{\linewidth}{1.1pt}}\vspace{1mm}

{\fontsize{7.4}{9}\selectfont\color{Slate}Ordered by cluster. Each entry gives the
methodological core and the finding that matters, not an abstract paraphrase.\par}

\band{Cardiovascular \& metabolic\hfill 1 record}{Teal}

\paperentry{De Potter et al.\ 2026}
{Socioeconomic and environmental determinants of STEMI incidence and in-hospital mortality}
{Eur Heart J Acute Cardiovasc Care}
{Nine years of Belgian national hospital discharge data, \textbf{51{,}544 STEMI
admissions} (51.4 per 100{,}000 per year), with residential PM$_{2.5}$, NO$_2$, O$_3$,
temperature, noise and greenspace from high-resolution models. A time-stratified
case-crossover handled the acute question and multivariable logistic regression the
mortality question. Per IQR, same-day: \textbf{PM$_{2.5}$ OR 1.022 (1.008--1.035)},
\textbf{NO$_2$ OR 1.044 (1.025--1.063)}; O$_3$ peaked over lag 0--2 at
\textbf{1.063 (1.034--1.090)}. In-hospital mortality was 11.8\% and tracked
\emph{socioeconomic} variables --- lower individual and area income, retirement status ---
rather than any pollutant. Two readings matter: NO$_2$ outperforming PM$_{2.5}$ in
single-pollutant models is the standard traffic-proximity confounding, and the
incidence/mortality split says exposure and outcome severity are governed by different
mechanisms.}
{10.1093/ehjacc/zuag118}

\band{Reproductive \& developmental\hfill 1 record}{Sky}

\paperentry{Zanini et al.\ 2026}
{Maternal air-pollution exposure and fetoplacental Doppler markers}
{Environ Pollut}
{Ships on metadata only. The Crossref deposit carries no abstract, the record is not in
PubMed or Europe PMC, and Semantic Scholar returns 404 for the DOI, so cohort size,
exposure model and Doppler effect sizes are all unavailable. Counted and listed; no
numbers claimed. Flagged for the abstract retry list --- fetoplacental haemodynamics is a
mechanistically informative intermediate between maternal PM exposure and birth outcome
and is thinly represented in this watch's corpus.}
{10.1016/j.envpol.2026.129046}

\band{Respiratory \& allergic\hfill 1 record}{Sage}

\paperentry{Li H et al.\ 2026}
{Multi-omics Mendelian randomisation of PM$_{2.5}$ and childhood asthma}
{ERJ Open Res}
{Two-sample MR on large-scale GWAS summary data gives \textbf{OR 1.897 (1.063--3.388),
$p$=0.030} for genetically instrumented PM$_{2.5}$ on childhood asthma. A TWAS identified
\textbf{70 genes} co-expressed with both, enriched in lysosome- and leukocyte-mediated
immunity, with \textbf{MEAF6} validated as protective and \textbf{RNF40} as a risk gene.
Two-step MR put \textbf{FUT10 as a positive mediator carrying 19.3\%} of the effect, with
CD200 and MANBA negative mediators; a phenome-wide MR flagged associations of these
targets with other diseases, which the authors correctly present as a caution against
treating them as drug targets. The limits are the usual MR limits and they are severe
here: the PM$_{2.5}$ instrument is a genetic proxy for a residentially correlated
environmental exposure, so horizontal pleiotropy through socioeconomic position is not
excluded, and the interval spans a factor of three.}
{10.1183/23120541.01444-2025}

\band{Occupational \& indoor\hfill 4 records}{Clay}

\paperentry{Lv H et al.\ 2026}
{Polyimide microsphere/nanofibre gradient aerogels for high-temperature filtration}
{ACS Appl Mater Interfaces}
{A single-component polyimide gradient aerogel built to beat the dust-holding limit of
uniform nanofibre mats, which cake at the surface and choke. Dual-layer structure:
$\approx$730\,nm coarse fibres under $\approx$300\,nm medium fibres, with a composite
functional layer of $\approx$3.34\,$\mu$m microspheres and $\approx$60\,nm ultrafine
fibres on top. Performance: \textbf{99.996\% PM$_{0.3}$ efficiency at 45.08\,Pa} initial
pressure drop, \textbf{127\,g\,m$^{-2}$ dust-holding capacity}, and \textbf{99.91\%
retained at 300\,\textdegree C}. Single-component construction is claimed for
recyclability and uniform thermal expansion. What is not reported is loaded pressure-drop
evolution or efficiency after thermal cycling, which is what actually determines service
life in a stack.}
{10.1021/acsami.6c10042}

\paperentry{Ma L et al.\ 2026}
{Surgical smoke from laser laparoscopic partial nephrectomy with an aspirating handle}
{Can J Urol}
{Ex vivo porcine kidney comparison of electrocoagulation, ultrasonic, 980\,nm diode and
532\,nm KTP laser resection, with and without a novel aspirator-integrated handle. The
KTP green laser produced the most smoke; coagulation-layer thickness ran diode
(\textbf{4.16 $\pm$ 0.21\,mm}) $>$ green (1.55 $\pm$ 0.21) $>$ electrocoagulation
(0.35 $\pm$ 0.08) $>$ ultrasonic (0.52 $\pm$ 0.08). With the aspirator, diode-laser
formaldehyde fell \textbf{4.95 $\to$ 1.69\,mg\,m$^{-3}$}, TVOC \textbf{8.34 $\to$
1.11\,mg\,m$^{-3}$} and CO \textbf{124 $\to$ 24\,ppm}. Gases only --- the study reports no
particle-size distribution or number concentration, which is the fraction that reaches the
alveoli and the reason surgical smoke is an occupational aerosol problem rather than a
gas one. In vitro, no personal sampling.}
{10.32604/cju.2026.070096}

\paperentry{Chen Y et al.\ 2026}
{Solid-fuel PAH exposure, phenanthrene metabolites, DNA damage and inflammation}
{Atmos Environ}
{Ships on metadata only --- no abstract in the Crossref deposit and no PubMed, Europe PMC
or Semantic Scholar record on this run. Title and authorship indicate urinary
phenanthrene metabolites as the exposure index in a rural northern Chinese solid-fuel
population, with DNA-damage and systemic-inflammation biomarkers as outcomes. Counted and
listed; no numbers claimed. High priority for the retry list: household solid-fuel
combustion is the single largest PM exposure contrast available in observational work and
this watch carries very little of it.}
{10.1016/j.atmosenv.2026.122339}

\paperentry{Micallef \& Scerri 2026}
{A flow-conditioned infiltration metric for traffic-related pollution ingress}
{Build Environ}
{Ships on metadata only --- Crossref carries no abstract and the record is not otherwise
indexed on this run. The title indicates an outdoor-to-indoor infiltration metric for
traffic-related air pollution in naturally ventilated buildings, conditioned on flow
regime rather than assumed constant. Counted and listed; no numbers claimed. Worth
retrying: a flow-dependent infiltration factor is exactly the term that makes
outdoor-monitor-based exposure assignment wrong for naturally ventilated stock, and it is
the term a paired indoor/outdoor sensor pair can actually estimate.}
{10.1016/j.buildenv.2026.115195}

\band{Sensing, forecasting \& instrumentation\hfill 2 records}{Amber}

\paperentry{Chaparro et al.\ 2026}
{Biological and technological sensors for airborne particles at heritage sites}
{Sci Total Environ}
{Ships on metadata only. The title indicates parallel biomonitoring and instrumental
sensing of airborne fine particles for preventive conservation of cultural heritage, but
no abstract was deposited and the record is not in PubMed, Europe PMC or Semantic Scholar
on this run, so the sensor models, the biomonitor species and any agreement statistics
between the two are unavailable. Counted and listed; no numbers claimed. On the retry
list --- cross-validation of a biological against an instrumental particle sensor is
directly relevant to low-cost network validation and rarely published.}
{10.1016/j.scitotenv.2026.182285}

\paperentry{Delkash 2026}
{Barometric pressure variability in community air monitoring near a landfill}
{Environ Pollut}
{Ships on metadata only --- the Crossref deposit has no abstract and Semantic Scholar
returns the record with a null abstract field. The title describes a receptor-scale
multi-method analysis of how barometric pressure variability imprints on community
monitoring data around an elevated-temperature landfill, across two mitigation
interventions. Counted and listed; no numbers claimed. This is the single most operationally
relevant record of the day for network work --- pressure-driven emission modulation is a
known confounder in fenceline and community monitoring and is almost never quantified ---
so it goes to the top of the retry list.}
{10.1016/j.envpol.2026.129084}

\band{Exposure assessment \& modelling\hfill 6 records}{Amber}

\paperentry{Petersen et al.\ 2026}
{Sulfate ions in aerosol water are their own source of radical oxidants}
{Environ Sci Technol}
{Direct spectroscopic evidence that the sulfate anion absorbs surface-level sunlight and
ejects an electron, forming \textbf{SO$_4^{\bullet-}$ on a sub-microsecond timescale} in a
concentration-dependent manner; long-timescale production rates quantified by radical
trapping with LC-HRMS. Compounds emitted from biomass fires photosensitise the pathway,
raising radical formation \textbf{up to five-fold} over direct photolysis across
pH 1.6--5.5, more strongly at low pH. Estimated SO$_4^{\bullet-}$ production is
\textbf{comparable to measured hydroxyl-radical production in aerosols} --- i.e.\ a second
oxidant budget nobody was counting. The downstream claim is the striking one: secondary
brown carbon from this chemistry \textbf{turns the direct radiative forcing of
sulfate-containing aerosol from negative to positive}, and the organosulfates formed are
cloud-active. Bulk solutions at controlled ionic strength; extrapolation to real
aerosol water is the open step.}
{10.1021/acs.est.6c05704}

\paperentry{Engeroff et al.\ 2026}
{PM exposure while running: park, residential and roadside routes by time of day}
{Research Square (preprint)}
{Mobile personal PM$_{10}$, PM$_{2.5}$ and PM$_1$ over \textbf{60 instrumented 2\,km runs}
(20 each on a major road, a residential street and an urban park; 30 morning, 30 midday;
pace 8:02--4:31\,min/km) with local urban background as reference and linear mixed-effects
models adjusting for meteorology. Lowest concentrations of every fraction:
\textbf{midday in the park}. Highest PM$_{10}$ anywhere: \textbf{morning on that same park
route}. Exposure across all routes was lower at midday, tracking urban background, and
humidity, pressure, wind speed and temperature all carried signal. The conclusion --- park
routes and rush-hour avoidance are sufficient --- is reasonable but rests on 20 runs per
route in one city, and the fact that the park produced both the minimum and the maximum
should temper any general route advice. Preprint, not peer reviewed.}
{10.21203/rs.3.rs-10016764/v1}

\paperentry{Zhu et al.\ 2026}
{Stabilized Criegee intermediates and the sensitivity of sulfuric acid formation}
{Atmos Chem Phys}
{Criegee reactions in MCM v3.3.1 updated to current IUPAC recommendations, driven through
a large AtChem box-model ensemble, with interpretable XGBoost surrogates used to attribute
the sCI-mediated SO$_2$ loss rate and its fractional share to controlling factors. The
sCI loss rate is governed chiefly by \textbf{relative humidity and isoprene}; the drivers
of the sCI \emph{fraction} split diurnally --- alkene speciation, O$_3$ and alkene share of
VOCs by day, alkene speciation, RH and NO$_x$ by night. Under high-sCI-fraction regimes,
the positive response of total SO$_2$ oxidation to O$_3$, VOCs and alkene fraction is
\emph{amplified}, which a summer episode at Wuhai supports. Policy consequence: where
alkene ozonolysis is active, \textbf{cutting VOCs controls sulfuric acid formation more
effectively than cutting SO$_2$ alone} --- and sulfuric acid is the new-particle-formation
precursor. Box model with a machine-learned surrogate; the surrogate inherits every
mechanistic assumption underneath it.}
{10.5194/acp-26-12479-2026}

\paperentry{Zhao H et al.\ 2026}
{Lead isotopes resolve dustfall sources across urban functional areas}
{Environ Monit Assess}
{Ten trace elements (V, Mn, Co, Cr, Ni, Cu, Zn, As, Cd, Pb) in dustfall from
\textbf{16 sites} across residential, busy-road, park and industrial areas of Baoding,
with enrichment factors, geo-accumulation index and Pb isotope fingerprinting. Cu, Pb, Zn
and Cd reach considerable pollution levels; the spatial ranking is
\textbf{residential $>$ busy road $>$ industrial $>$ park}, which is not the ordering most
source-apportionment intuition predicts and the paper does not fully explain. Total
non-carcinogenic risk stayed below threshold but \textbf{total carcinogenic risk exceeded
the acceptable threshold} in several areas, worst on busy roads and in residential zones.
Pb isotopes attribute \textbf{over 50\%} to a combined coal-combustion plus
traffic-emission group in most areas --- combined, i.e.\ the isotope system does not
separate the two here, which is the honest limit of the method in a city that burns coal
and drives cars.}
{10.1007/s10661-026-15859-0}

\paperentry{Lei et al.\ 2026}
{Source transition in PM$_{2.5}$-bound metals in a coastal megacity}
{Atmos Environ}
{Ships on metadata only --- no abstract in the Crossref deposit and no other index carries
the record on this run. The title indicates a multi-year, high-time-resolution
observational record of PM$_{2.5}$-bound metals with a source-transition analysis and a
health-implication component. Counted and listed; no numbers claimed. High retry priority:
multi-year hourly metals is exactly the dataset class against which a low-cost network's
composition proxies would have to be validated, and there are very few of them.}
{10.1016/j.atmosenv.2026.122336}

\paperentry{Mutic et al.\ 2026}
{Paired silicone wristbands across home and early-childhood-education settings}
{medRxiv (preprint)}
{A \emph{protocol}, not results. Prospective repeated-measures pilot targeting
\textbf{44 preschool children and 8 staff} across two sociodemographically distinct Atlanta
early-childhood centres, with home- and centre-designated silicone wristbands exchanged
between settings over three consecutive days and nights, one spot urine per child, and
continuous indoor air-quality monitoring in two classrooms per site. Eight pre-specified
feasibility indicators. The design premise is sound and underexploited --- children spend
$>$90\% of time indoors and early-childhood centres are high-occupant-density
microenvironments almost absent from the exposure literature --- but wristbands integrate
semi-volatile organics, not particle mass, so the PM component rests entirely on the
fixed classroom monitors. Listed for the design; nothing to conclude from it yet.}
{10.64898/2026.08.27.26361470}

\clearpage

% ==================================================================== TABLE
\section*{\sffamily\bfseries\color{Deep}Corpus register}
\vspace{-1.5mm}{\color{Amber}\rule{\linewidth}{1.1pt}}\vspace{2.5mm}

{\fontsize{6.7}{8.4}\selectfont
\setlength{\tabcolsep}{2.4pt}
\begin{longtable}{@{}>{\RaggedRight}p{25mm}>{\RaggedRight}p{41mm}>{\RaggedRight}p{27mm}%
>{\RaggedRight}p{23mm}>{\RaggedRight}p{18mm}>{\RaggedRight}p{15mm}@{}}
\toprule
\textsf{\textbf{First author}} & \textsf{\textbf{Focus}} & \textsf{\textbf{Journal}} &
\textsf{\textbf{Design}} & \textsf{\textbf{Metric}} & \textsf{\textbf{Region}}\\
\midrule
\endfirsthead
\toprule
\textsf{\textbf{First author}} & \textsf{\textbf{Focus}} & \textsf{\textbf{Journal}} &
\textsf{\textbf{Design}} & \textsf{\textbf{Metric}} & \textsf{\textbf{Region}}\\
\midrule
\endhead
\multicolumn{6}{@{}l@{}}{\cellcolor{Mist}\textsf{\textbf{\textcolor{Deep}{Neuro / mental health}}}}\\[0.6mm]
\href{https://doi.org/10.1001/jamanetworkopen.2026.25015}{Dales} & Wildfire-sourced PM$_{2.5}$ and ED visits for migraine and other primary headache & \textit{JAMA Netw Open} & Case-crossover & PM$_{2.5}$ (wildfire) & Canada\\
\href{https://doi.org/10.1093/ageing/afag215}{Glover} & Umbrella meta-review of urban environmental exposures and cognitive health & \textit{Age Ageing} & Umbrella review / meta-analysis & PM$_{2.5}$ & Global\\
\href{https://doi.org/10.1016/j.nsa.2026.107019}{Goossens} & In utero air pollution exposure and pre-/postnatal brain development & \textit{Neurosci Appl} & Narrative review & PM$_{2.5}$ / PM$_{10}$ / PM$_{\mathrm{coarse}}$ & Global\\
\href{https://doi.org/10.1016/j.isci.2026.116617}{Ma} & PM$_{2.5}$-induced neurotoxicity and depression: exposure risks and mechanisms & \textit{iScience} & Narrative review & PM$_{2.5}$ + components & Global\\
\href{https://doi.org/10.3760/cma.j.cn112150-20250804-00752}{Pan} & Maternal PM exposure and offspring neurodevelopment: effects and mechanisms & \textit{Zhonghua Yu Fang Yi Xue Za Zhi} & Narrative review & PM (bulk) & China\\
\addlinespace[0.9mm]
\multicolumn{6}{@{}l@{}}{\cellcolor{Mist}\textsf{\textbf{\textcolor{Teal}{Cardiovascular \& metabolic}}}}\\[0.6mm]
\href{https://doi.org/10.3760/cma.j.cn112338-20251120-00843}{Huang} & PM$_{2.5}$ and dyslipidaemia in rural Chinese children; dietary effect modification & \textit{Zhonghua Liu Xing Bing Xue Za Zhi} & Cross-sectional (multistage) & PM$_{2.5}$ & China\\
\href{https://doi.org/10.32394/pe/218966}{Kazmierska} & Air pollution and cardiovascular disease: review of current data & \textit{Przegl Epidemiol} & Narrative review & PM$_{2.5}$ / PM$_{10}$ & Global\\
\href{https://doi.org/10.1007/s00125-026-06814-2}{Xi} & Air pollution, built environment and temperature vs. prediabetes phenotypes & \textit{Diabetologia} & Prospective cohort & PM$_{2.5}$ & Germany\\
\href{https://doi.org/10.1161/JAHA.125.046917}{Yang} & Short-term PM$_{2.5}$ and hospitalisation/ED use in atrial fibrillation & \textit{J Am Heart Assoc} & Case-crossover & PM$_{2.5}$ (wildfire-influenced) & USA\\
\addlinespace[0.9mm]
\multicolumn{6}{@{}l@{}}{\cellcolor{Mist}\textsf{\textbf{\textcolor{Sky}{Reproductive \& developmental}}}}\\[0.6mm]
\href{https://doi.org/10.1016/j.jes.2025.10.043}{Chen} & Prenatal air pollution, fetal growth trajectory phenotypes and child neurodevelopment & \textit{J Environ Sci (China)} & Prospective cohort & PM$_{2.5}$ / PM$_{10}$ & China\\
\href{https://doi.org/10.1016/j.ecoenv.2026.120566}{Huang} & Joint prenatal bisphenol + air-pollutant mixture exposure and preschool IQ & \textit{Ecotoxicol Environ Saf} & Prospective cohort (mixtures) & PM$_{2.5}$ / PM$_{10}$ & China\\
\href{https://doi.org/10.1007/s40572-026-00552-8}{Peterson} & Ambient air pollution, PAHs and pubertal development: critical review & \textit{Curr Environ Health Rep} & Critical review & PM$_{2.5}$ components & Global\\
\href{https://doi.org/10.1016/j.envres.2026.125313}{Xie} & Long-term PM$_{2.5}$ and timing of pubertal development in boys and girls & \textit{Environ Res} & Prospective cohort & PM$_{2.5}$ & China\\
\addlinespace[0.9mm]
\multicolumn{6}{@{}l@{}}{\cellcolor{Mist}\textsf{\textbf{\textcolor{Sage}{Respiratory \& allergic}}}}\\[0.6mm]
\href{https://doi.org/10.1007/s00484-026-03289-z}{Bayar Muluk} & Climate change and otolaryngology: airway disease impacts and emission reduction & \textit{Int J Biometeorol} & Narrative review & PM$_{2.5}$ & Global\\
\href{https://doi.org/10.4103/mgr.MEDGASRES-D-25-00164}{Liao} & Short-term air pollution, respiratory symptoms and lung function in asthma & \textit{Med Gas Res} & Panel study & PM$_{2.5}$ / PM$_{10}$ & China\\
\href{https://doi.org/10.3390/epidemiologia7040097}{Sanchez de Prada} & Ambient PM and asthma-coded admissions in children 0-3 y, Spain & \textit{Epidemiologia} & Ecological panel & PM$_{2.5}$ / PM$_{10}$ & Spain\\
\href{https://doi.org/10.1016/j.jhazmat.2026.143076}{Yang} & Allergenic vs. toxic risk of airborne seed hairs compared with ambient PM$_{2.5}$ & \textit{J Hazard Mater} & Field measurement & PM$_{2.5}$ + bioaerosol & China\\
\addlinespace[0.9mm]
\multicolumn{6}{@{}l@{}}{\cellcolor{Mist}\textsf{\textbf{\textcolor{Violet}{Mechanistic toxicology}}}}\\[0.6mm]
\href{https://doi.org/10.1016/j.jes.2025.12.074}{Ding} & PM$_{2.5}$ promotes vascular calcification via miR-27b-5p/HES1 signalling & \textit{J Environ Sci (China)} & In vitro + organoid + murine & PM$_{2.5}$ & China\\
\href{https://doi.org/10.1007/s00360-026-01698-5}{Maraschi} & Osmoregulatory disruption and metal bioaccumulation from settleable atmospheric PM & \textit{J Comp Physiol B} & Controlled exposure (ecotox) & Settleable APM & Brazil\\
\href{https://doi.org/10.1016/j.jes.2025.11.018}{Wu} & Macrophage NLRP3 inflammasome in real-ambient PM-induced hepatic glucose dysregulation & \textit{J Environ Sci (China)} & Real-ambient murine + co-culture & Real-ambient PM & China\\
\addlinespace[0.9mm]
\multicolumn{6}{@{}l@{}}{\cellcolor{Mist}\textsf{\textbf{\textcolor{Clay}{Occupational \& indoor}}}}\\[0.6mm]
\href{https://doi.org/10.3389/fmedt.2026.1796828}{Chen} & Nested suction instrument to cut intraoperative surgical-smoke PM$_{2.5}$ & \textit{Front Med Technol} & Bench / device experiment & PM$_{2.5}$ (surgical smoke) & China\\
\href{https://doi.org/10.1093/annweh/wxag055}{Cole} & Determinants of respirable crystalline silica exposure in Australian tunnelling & \textit{Ann Work Expo Health} & Observational exposure study & Respirable dust / RCS & Australia\\
\href{https://doi.org/10.1016/j.jhazmat.2026.143081}{Ding} & Stage-resolved hair Zn-isotope dynamics as a firefighting combustion-exposure tracer & \textit{J Hazard Mater} & Repeated-measures biomarker & Combustion-derived PM & China\\
\href{https://doi.org/10.3390/jox16040119}{Ramsden} & Xenobiotic hazards in aircraft cabin air & \textit{J Xenobiot} & Review + risk economics & Ultrafine particles & Global\\
\addlinespace[0.9mm]
\multicolumn{6}{@{}l@{}}{\cellcolor{Mist}\textsf{\textbf{\textcolor{Amber}{Sensing, forecasting \& instrumentation}}}}\\[0.6mm]
\href{https://doi.org/10.1016/j.jenvman.2026.130515}{Shi} & Adaptive spatiotemporal graph transformer for multi-site multi-step PM$_{2.5}$ forecasting & \textit{J Environ Manage} & Deep-learning model & PM$_{2.5}$ & China / India\\
\href{https://doi.org/10.3390/bios16070380}{Tsuchiya} & Seasonal and indoor-outdoor characteristics of fluorescent aerosol particles in offices & \textit{Biosensors} & Real-time sensor campaign & Bioaerosol / coarse & Japan\\
\addlinespace[0.9mm]
\multicolumn{6}{@{}l@{}}{\cellcolor{Mist}\textsf{\textbf{\textcolor{Amber}{Exposure assessment \& modelling}}}}\\[0.6mm]
\href{https://doi.org/10.1029/2025GH001621}{Solorzano-Araque} & Multiscale exposure modelling: LOTOS-EUROS + mobility surveys + particle toxicity & \textit{GeoHealth} & CTM + mobility integration & PM$_{2.5}$ / PM$_{10}$ & Colombia\\
\addlinespace[0.9mm]
\multicolumn{6}{@{}l@{}}{\cellcolor{Mist}\textsf{\textbf{\textcolor{Coral}{Burden, policy \& mitigation}}}}\\[0.6mm]
\href{https://doi.org/10.1002/hsr2.72886}{Karamzad} & Ischaemic heart disease burden in MENA 1990-2021; PM attributable fraction 28.2\% & \textit{Health Sci Rep} & GBD secondary analysis & Ambient PM & MENA\\
\href{https://doi.org/10.1371/journal.pone.0347989}{Rees} & HEPA filtration to reduce respiratory infection in care homes & \textit{PLoS One} & RCT process evaluation & Filtered PM & UK\\
\href{https://doi.org/10.1016/j.jes.2026.01.009}{Shan} & Policy-driven emission reductions and post-pandemic rebound, high-resolution inventories & \textit{J Environ Sci (China)} & Emission inventory & PM emissions & China\\
\addlinespace[0.9mm]
\multicolumn{6}{@{}l@{}}{\cellcolor{Mist}\textsf{\textbf{\textcolor{Slate}{Other clinical endpoints}}}}\\[0.6mm]
\href{https://doi.org/10.3390/dj14070440}{Lupita} & PM$_{2.5}$, preventive-care capacity and caries/edentulism burden across the EU & \textit{Dent J} & Ecological panel & PM$_{2.5}$ & EU-27\\
\href{https://doi.org/10.1016/j.jdin.2026.06.008}{Park} & Long-term PM exposure and risk of androgenetic alopecia & \textit{JAAD Int} & Nationwide cohort & PM$_{2.5}$ / PM$_{10}$ & South Korea\\
\href{https://doi.org/10.3171/2026.3.PEDS2615}{Sun} & Environmental exposure and area deprivation vs. paediatric hydrocephalus outcomes & \textit{J Neurosurg Pediatr} & Registry cohort + geospatial & PM$_{2.5}$ (satellite-derived) & USA\\
\addlinespace[0.9mm]
\bottomrule
\end{longtable}}


\vspace{4mm}
\begin{keybox}
{\sffamily\bfseries\footnotesize\color{Deep}Provenance and method}\\[1.2mm]
{\fontsize{7.2}{8.8}\selectfont
Entry window \textbf{2 September 2026}, single day, opened against a
\texttt{last\_entry\_date} of 1 September. Sources queried on 3 September 2026; under this
pipeline's 22:00-local rule, 2 September is the newest buildable date and 3 September is
not opened.
\textbf{PubMed} E-utilities on \texttt{[EDAT]} 2026/09/02: 4 hits on the health axis and
\textbf{0} on the instrumentation axis.
\textbf{Europe PMC} \texttt{CREATION\_DATE}: 17 hits, including the two preprints carried
here.
\textbf{Crossref} by-ISSN, windowed on \texttt{created}: 68 deposits.
\textbf{OpenAlex} HTTP 403 (plan-gated date filter) and \textbf{arXiv} HTTP 429; both legs
recorded as failed, not empty. Consensus was not re-queried --- it returned three
already-seen records on the 1 September run and its calibration axis does not move daily.
After collapsing cross-source duplicates by DOI and normalised title, 89 unique records
were considered; \textbf{12 were already shipped} (11 on 1 September, 1 on 31 August ---
the by-ISSN sweep re-lists a deposit for several days), 15 were screened in scope and 62
were screened out. Notable exclusions: an \textit{Atmos.\ Environ.} O$_3$-only cohort
(\texttt{10.1016/j.atmosenv.2026.122338}) and a Barcelona indoor-NO$_2$ modelling study,
both excluded as non-particulate; a settled-dust contaminants-of-emerging-concern survey,
excluded as a chemical rather than aerosol measurement.
\textbf{Six of the fifteen in-scope records ship on metadata only} and carry no attributed
numbers; every one is on the abstract retry list.
Subtopic, design, particle-metric and geography labels are single-label by dominant
contribution and are a judgement, not a controlled vocabulary. All DOI links resolve to
the publisher of record; abstracts remain the copyright of their respective publishers.\par}
\end{keybox}

\end{document}
